%% ----------------------------------------------------------------------------
%%  §11.5  La serie armonica
%% ----------------------------------------------------------------------------
\section{La serie armonica}
\label{sec:serie-armonica}

La serie
\begin{equation}\label{eq:serie-armonica}
  \sum_{k=1}^{\infty}\frac1k=1+\frac12+\frac13+\dots+\frac1n+\cdots
\end{equation}
si chiama \emph{serie armonica}\index{serie!armonica}. Il suo termine generale tende a zero,
e le somme parziali crescono molto lentamente: ci vogliono più di dodicimila termini
perché superino $10$. Eppure la serie diverge. Poiché il risultato è importante e un po'
sorprendente, ne diamo tre dimostrazioni, ciascuna delle quali insegna qualcosa.

\begin{teorema}[Divergenza della serie armonica]\label{teo:armonica}
La serie armonica è divergente: $\sum_{k=1}^{\infty}\frac1k=+\infty$. Più precisamente, le sue
somme parziali $H_{n}=\sum_{k=1}^{n}\frac1k$ verificano
\begin{equation}\label{eq:armonica-log}
  H_{n}\ge\log(n+1) \qquad\perogni n\in\N .
\end{equation}
\end{teorema}

\begin{schema}
\Ipotesi{$H_{n}=\sum_{k=1}^{n}\frac1k$.}
\Tesi{$H_{n}\ge\log(n+1)$ per ogni $n$; \ quindi $H_{n}\to+\infty$.}
\end{schema}

\begin{strategia}
Confrontiamo la somma con un'area. Il termine $\frac1k$ è l'area di un rettangolo di base
$[k,k+1]$ e altezza $\frac1k$, che contiene la regione sotto il grafico di $\frac1x$ nello stesso
intervallo, perché $\frac1x\le\frac1k$ per $x\ge k$ (\cref{fig:armonica-integrale}, a sinistra).
Sommando, $H_{n}$ supera l'integrale di $\frac1x$ tra $1$ e $n+1$, cioè $\log(n+1)$, che diverge.
\end{strategia}

\begin{dimostrazione}[Prima dimostrazione (con l'integrale)]
Per ogni $k\in\N$ e ogni $x\in[k,k+1]$ si ha $\frac1x\le\frac1k$. Per il confronto tra integrali
(\cref{teo:confronto-integrali}),
\begin{equation}\label{eq:armonica-rettangolo}
  \int_{k}^{k+1}\frac{dx}x\le\int_{k}^{k+1}\frac{dx}k=\frac1k .
\end{equation}
Sommando per $k$ da $1$ a $n$ e usando l'additività dell'integrale (\cref{teo:additivita})
e la formula fondamentale del calcolo integrale,
\[
  H_{n}=\sum_{k=1}^{n}\frac1k\ge\sum_{k=1}^{n}\int_{k}^{k+1}\frac{dx}x=\int_{1}^{n+1}\frac{dx}x=\log(n+1) .
\]
Poiché $\log(n+1)\to+\infty$, per il teorema del confronto (\cref{teo:confronto-infinito})
anche $H_{n}\to+\infty$.
\end{dimostrazione}

La seconda dimostrazione ottiene la stessa disuguaglianza senza il calcolo integrale,
usando le proprietà del numero $e$.

\begin{dimostrazione}[Seconda dimostrazione (con il numero $e$)]
Per il \cref{teo:numero-e} si ha $\bigl(1+\frac1k\bigr)^{k}<e$ per ogni $k\in\N$. Poiché il
logaritmo è crescente,
\[
  k\log\Bigl(1+\frac1k\Bigr)=\log\Bigl(1+\frac1k\Bigr)^{k}<\log e=1 ,
\]
cioè
\[
  \frac1k>\log\frac{k+1}k=\log(k+1)-\log k .
\]
Sommando per $k=1,\dots,n$, il secondo membro è una somma telescopica
(\cref{prop:serie-telescopica}):
$H_{n}>\log(n+1)-\log1=\log(n+1)$.
\end{dimostrazione}

La terza dimostrazione, la più antica, raggruppa i termini a blocchi e usa il criterio di
Cauchy.

\begin{dimostrazione}[Terza dimostrazione (con il criterio di Cauchy)]
Per ogni $n\in\N$, il blocco dei $n$ termini da $\frac1{n+1}$ a $\frac1{2n}$ verifica
\[
  \sum_{k=n+1}^{2n}\frac1k=\frac1{n+1}+\frac1{n+2}+\dots+\frac1{2n}\ge n\cdot\frac1{2n}=\frac12 ,
\]
perché ciascuno dei suoi $n$ termini è almeno $\frac1{2n}$. Dunque la condizione
\eqref{eq:cauchy-serie} è violata con $\epsilon=\frac12$ e $p=n$, per ogni $n$: la serie non
converge (\cref{teo:cauchy-serie}). Essendo a termini positivi, diverge a $+\infty$
(\cref{teo:serie-positive}).
\end{dimostrazione}

\begin{figure}[htb]
\centering
\begin{tikzpicture}
\begin{axis}[width=0.53\textwidth, height=4.6cm, xmin=0, xmax=8.6, ymin=0, ymax=1.15,
  axis lines=left, xtick={1,2,...,8}, ytick={0,0.5,1}, axis on top,
  title={\scriptsize $H_{n}\ge\int_{1}^{n+1}\frac{dx}x$}, tick label style={font=\scriptsize}]
  \addplot[fill=blu!20, draw=blu, ybar interval]
    coordinates {(1,1) (2,0.5) (3,0.3333) (4,0.25) (5,0.2) (6,0.1667) (7,0.1429) (8,0)};
  \addplot[curva, domain=0.87:8.6, samples=80] {1/x};
\end{axis}
\end{tikzpicture}\hfill
\begin{tikzpicture}
\begin{axis}[width=0.53\textwidth, height=4.6cm, xmin=0, xmax=8.6, ymin=0, ymax=1.15,
  axis lines=left, xtick={1,2,...,8}, ytick={0,0.5,1}, axis on top,
  title={\scriptsize $H_{n+1}-1\le\int_{1}^{n+1}\frac{dx}x$}, tick label style={font=\scriptsize}]
  \addplot[fill=rame!30, draw=rame, ybar interval]
    coordinates {(1,0.5) (2,0.3333) (3,0.25) (4,0.2) (5,0.1667) (6,0.1429) (7,0.125) (8,0)};
  \addplot[curva, domain=0.87:8.6, samples=80] {1/x};
\end{axis}
\end{tikzpicture}
\caption{Il confronto tra la serie armonica e l'integrale di $\frac1x$. A sinistra i
rettangoli di altezza $\frac1k$ sulle basi $[k,k+1]$ contengono la regione sotto il grafico:
$H_{n}\ge\log(n+1)$. A destra i rettangoli di altezza $\frac1{k+1}$ sono contenuti in essa:
$H_{n+1}-1\le\log(n+1)$. Le somme parziali sono quindi «intrappolate» tra due logaritmi.}
\label{fig:armonica-integrale}
\end{figure}

\subsection{Il criterio dell'integrale e la serie armonica generalizzata}

Il ragionamento della prima dimostrazione funziona per ogni funzione positiva e
decrescente, e collega il carattere di una serie a quello di un integrale improprio
(\S\ref{sec:impropri}).

\begin{teorema}[Criterio dell'integrale]\label{teo:criterio-integrale}
\index{criterio dell'integrale}
Sia $f$ continua, non negativa e decrescente in $[1,+\infty)$, e sia $a_{k}=f(k)$. Allora
per ogni $n\in\N$
\begin{equation}\label{eq:criterio-integrale}
  \int_{1}^{n+1}f(x)\,dx\le\sum_{k=1}^{n}a_{k}\le a_{1}+\int_{1}^{n}f(x)\,dx ,
\end{equation}
e la serie $\sum_{k=1}^{\infty}a_{k}$ converge se e solo se converge l'integrale improprio
$\int_{1}^{+\infty}f(x)\,dx$. In tal caso il resto $R_{n}=\sum_{k=n+1}^{\infty}a_{k}$ verifica
\begin{equation}\label{eq:resto-integrale-serie}
  \int_{n+1}^{+\infty}f(x)\,dx\le R_{n}\le\int_{n}^{+\infty}f(x)\,dx .
\end{equation}
\end{teorema}

\begin{schema}
\Ipotesi{$f$ continua, $f\ge0$, decrescente in $[1,+\infty)$; \ $a_{k}=f(k)$.}
\Tesi{la \eqref{eq:criterio-integrale}; \ $\sum a_{k}<+\infty$ $\sse$
$\int_{1}^{+\infty}f<+\infty$; \ la stima \eqref{eq:resto-integrale-serie} del resto.}
\end{schema}

\begin{strategia}
Come nella \cref{fig:armonica-integrale}: nell'intervallo $[k,k+1]$ la funzione
decrescente sta tra $f(k+1)$ e $f(k)$, e integrando si ottiene
$a_{k+1}\le\int_{k}^{k+1}f\le a_{k}$. Sommando si ottengono le due stime. Le somme parziali e la
funzione integrale $b\mapsto\int_{1}^{b}f$ sono entrambe crescenti, e le stime dicono che una è
limitata se e solo se lo è l'altra.
\end{strategia}

\begin{dimostrazione}
Per $x\in[k,k+1]$ la monotonia dà $f(k+1)\le f(x)\le f(k)$, e per il confronto tra
integrali
\begin{equation}\label{eq:integrale-intrappolato}
  a_{k+1}\le\int_{k}^{k+1}f(x)\,dx\le a_{k} \qquad\perogni k\in\N .
\end{equation}
Sommando la disuguaglianza di destra per $k=1,\dots,n$ si ottiene
$\int_{1}^{n+1}f\le\sum_{k=1}^{n}a_{k}$; sommando quella di sinistra per $k=1,\dots,n-1$ si ottiene
$\sum_{k=2}^{n}a_{k}\le\int_{1}^{n}f$, cioè la seconda delle \eqref{eq:criterio-integrale}.

Poiché $f\ge0$, la funzione $F(b)=\int_{1}^{b}f$ è crescente, e per il \cref{lem:monotona-infinito}
ha limite per $b\to+\infty$, finito se e solo se $F$ è limitata; inoltre, essendo crescente,
$F$ è limitata se e solo se lo è la successione $F(n)$. Le \eqref{eq:criterio-integrale}
mostrano che $s_{n}$ è limitata se e solo se $F(n)$ lo è, e per il \cref{teo:serie-positive}
la serie converge se e solo se converge l'integrale.

Infine, sommando le \eqref{eq:integrale-intrappolato} per $k$ da $n+1$ a $m$ si ha
$\int_{n+1}^{m+1}f\le\sum_{k=n+1}^{m}a_{k}$ e, per $k$ da $n$ a $m-1$, $\sum_{k=n+1}^{m}a_{k}\le\int_{n}^{m}f$;
per $m\to+\infty$ si ottiene la \eqref{eq:resto-integrale-serie}.
\end{dimostrazione}

\begin{teorema}[Serie armonica generalizzata]\label{teo:armonica-generalizzata}
\index{serie!armonica generalizzata}
La serie
\begin{equation}\label{eq:armonica-generalizzata}
  \sum_{k=1}^{\infty}\frac1{k^{p}}=1+\frac1{2^{p}}+\frac1{3^{p}}+\dots+\frac1{n^{p}}+\cdots
\end{equation}
è convergente se $p>1$ ed è divergente se $p\le1$.
\end{teorema}

\begin{schema}
\Ipotesi{$p\in\R$.}
\Tesi{$\sum\frac1{k^{p}}<+\infty$ $\sse$ $p>1$.}
\end{schema}

\begin{strategia}
Per $p>0$ la funzione $\frac1{x^{p}}$ è positiva e decrescente, e il criterio dell'integrale
riduce il problema all'integrale improprio $\int_{1}^{+\infty}\frac{dx}{x^{p}}$, già calcolato nel
\cref{cap:integrali-indefiniti}. Per $p\le0$ il termine generale non tende a zero.
\end{strategia}

\begin{dimostrazione}
Se $p\le0$, $\frac1{k^{p}}=k^{-p}\ge1$ non tende a zero, e la serie (a termini positivi) diverge
per il \cref{teo:condizione-necessaria}. Se $p>0$, la funzione $f(x)=x^{-p}$ è continua,
positiva e decrescente in $[1,+\infty)$; per il \cref{teo:criterio-integrale} e la
\eqref{eq:impropri-p-infinito}, la serie converge se e solo se $p>1$.
\end{dimostrazione}

Per $p>1$ la \eqref{eq:criterio-integrale} dà anche
$\frac{(n+1)^{1-p}-1}{1-p}\le s_{n}\le1+\frac{1-n^{1-p}}{p-1}$, e passando al limite
$\frac1{p-1}\le\sum_{k=1}^{\infty}\frac1{k^{p}}\le\frac p{p-1}$. Ad esempio $1\le\sum\frac1{k^{2}}\le2$; il valore esatto
è $\frac{\pi^{2}}6\approx1{,}6449$, come scoprì Eulero nel 1734. La stima del resto
\eqref{eq:resto-integrale-serie} con $f(x)=\frac1{x^{2}}$ dice che
$\frac1{n+1}\le\frac{\pi^{2}}6-\sum_{k=1}^{n}\frac1{k^{2}}\le\frac1n$: con cento termini si ottiene una sola cifra
decimale esatta (l'errore è circa $10^{-2}$).

\subsection{La costante di Eulero--Mascheroni}

Le somme parziali della serie armonica divergono, ma in modo molto regolare: sono vicine
a $\log n$, e la differenza ha un limite.

\begin{proposizione}[Costante di Eulero--Mascheroni]\label{prop:eulero-mascheroni}
\index{costante di Eulero--Mascheroni}
La successione $d_{n}=H_{n}-\log n$ è decrescente e verifica $0\le d_{n}\le1$. Il suo limite
\begin{equation}\label{eq:eulero-mascheroni}
  \gamma=\lim_{n\to+\infty}\Bigl(1+\frac12+\dots+\frac1n-\log n\Bigr)=0{,}5772156649\ldots
\end{equation}
si chiama \emph{costante di Eulero--Mascheroni}. In particolare $H_{n}=\log n+\gamma+o(1)$.
\end{proposizione}

\begin{schema}
\Ipotesi{$d_{n}=H_{n}-\log n$.}
\Tesi{$d_{n+1}\le d_{n}$; \ $0\le d_{n}\le1$; \ quindi $d_{n}\to\gamma\in[0,1]$.}
\end{schema}

\begin{strategia}
Basta mostrare che $d_{n}$ è monotòna e limitata. La differenza $d_{n}-d_{n+1}$ è l'area compresa,
nell'intervallo $[n,n+1]$, tra il grafico di $\frac1x$ e il rettangolo di altezza $\frac1{n+1}$:
è positiva. La limitazione dal basso viene dalla stima \eqref{eq:armonica-log}, che dice
che $H_{n}$ supera $\log(n+1)$, e quindi anche $\log n$.
\end{strategia}

\begin{dimostrazione}
Si ha $d_{n}-d_{n+1}=\log(n+1)-\log n-\frac1{n+1}=\int_{n}^{n+1}\bigl(\frac1x-\frac1{n+1}\bigr)dx\ge0$,
perché l'integrando è non negativo in $[n,n+1]$: la successione è decrescente, e quindi
$d_{n}\le d_{1}=1$. Inoltre, per la \eqref{eq:armonica-log},
$d_{n}\ge\log(n+1)-\log n>0$. Per il teorema sulle successioni monotòne $d_{n}$ converge a un
limite $\gamma\in[0,1]$.
\end{dimostrazione}

La relazione $H_{n}\approx\log n+\gamma$ spiega la lentezza della divergenza: per avere
$H_{n}>10$ occorre $n\approx e^{10-\gamma}\approx12\,367$, e per $H_{n}>100$ un numero di termini
dell'ordine di $10^{43}$. Non si sa ancora se $\gamma$ sia un numero razionale.

\begin{python}[Serie armonica, costante di Eulero e problema di Basilea]
import math

H, n = 0.0, 0
while H <= 10:                        # primo n con H_n > 10
    n += 1
    H += 1 / n
print(f"H_n > 10 per la prima volta con n = {n}")

H, N = 0.0, 1
for k in range(1, 10**6 + 1):
    H += 1 / k
    if k == N:
        print(f"n = {k:>7}:  H_n - log n = {H - math.log(k):.8f}")
        N *= 10

s = sum(1 / k**2 for k in range(1, 101))
print(f"pi^2/6 - s_100 = {math.pi**2 / 6 - s:.6f}"
      f"   (stima: tra {1/101:.6f} e {1/100:.6f})")
\end{python}
\begin{risultato}
H_n > 10 per la prima volta con n = 12367
n =       1:  H_n - log n = 1.00000000
n =      10:  H_n - log n = 0.62638316
n =     100:  H_n - log n = 0.58220733
n =    1000:  H_n - log n = 0.57771558
n =   10000:  H_n - log n = 0.57726566
n =  100000:  H_n - log n = 0.57722066
n = 1000000:  H_n - log n = 0.57721616
pi^2/6 - s_100 = 0.009950   (stima: tra 0.009901 e 0.010000)
\end{risultato}

\subsection{Approfondimento: il criterio di condensazione di Cauchy}
\label{subsec:condensazione}

La terza dimostrazione della divergenza della serie armonica raggruppava i termini in
blocchi sempre più lunghi. Se i blocchi hanno lunghezza $1,2,4,8,\dots$, l'idea funziona
per ogni successione non negativa e decrescente, e conduce a un criterio, dovuto a
Cauchy (\emph{Cours d'analyse}, 1821), che non richiede il calcolo integrale: la serie
data ha lo stesso carattere di una serie formata soltanto dai termini di indice $2^{k}$,
ciascuno moltiplicato per $2^{k}$, cioè per il numero dei termini del blocco che
rappresenta.

\begin{teorema}[Criterio di condensazione]\label{teo:condensazione}
\index{criterio di condensazione}
Sia $a_{n}$ una successione non negativa e decrescente. Allora le serie
\begin{equation}\label{eq:condensazione}
  \sum_{n=1}^{\infty}a_{n} \qquad\text{e}\qquad
  \sum_{k=0}^{\infty}2^{k}a_{2^{k}}=a_{1}+2a_{2}+4a_{4}+8a_{8}+\cdots
\end{equation}
hanno lo stesso carattere. Più precisamente, dette $s_{n}=\sum_{j=1}^{n}a_{j}$ e
$t_{m}=\sum_{k=0}^{m}2^{k}a_{2^{k}}$ le somme parziali delle due serie, per ogni intero $m\ge0$
\begin{equation}\label{eq:condensazione-stime}
  \tfrac12\,t_{m}\le s_{2^{m}} \qquad\text{e}\qquad s_{2^{m+1}-1}\le t_{m} .
\end{equation}
\end{teorema}

\begin{schema}
\Ipotesi{$a_{n}\ge0$; \ $a_{n+1}\le a_{n}$ per ogni $n\in\N$.}
\Tesi{$\sum a_{n}<+\infty$ $\sse$ $\sum 2^{k}a_{2^{k}}<+\infty$; \ le stime
\eqref{eq:condensazione-stime}.}
\end{schema}

\begin{strategia}
Si dividono gli indici in blocchi compresi tra due potenze di $2$ consecutive. Il blocco
da $2^{k}$ a $2^{k+1}-1$ ha $2^{k}$ termini, tutti non maggiori del primo, $a_{2^{k}}$, perché la
successione decresce; il blocco da $2^{k-1}+1$ a $2^{k}$ ha $2^{k-1}$ termini, tutti non minori
dell'ultimo, $a_{2^{k}}$. Sommando sui blocchi si ottengono le due stime, e per le serie a
termini non negativi la limitatezza delle somme parziali passa dall'una all'altra.
\end{strategia}

\begin{dimostrazione}
Per ogni $k\ge0$ il blocco degli indici $j$ con $2^{k}\le j\le2^{k+1}-1$ contiene $2^{k}$ indici,
e per la monotonia $a_{j}\le a_{2^{k}}$; quindi la somma dei termini del blocco è al più
$2^{k}a_{2^{k}}$. Gli indici da $1$ a $2^{m+1}-1$ sono l'unione dei blocchi con $k=0,1,\dots,m$, e
sommando si ottiene $s_{2^{m+1}-1}\le\sum_{k=0}^{m}2^{k}a_{2^{k}}=t_{m}$.

Per ogni $k\ge1$ il blocco degli indici $j$ con $2^{k-1}+1\le j\le2^{k}$ contiene $2^{k-1}$ indici, e
$a_{j}\ge a_{2^{k}}$; la somma dei suoi termini è quindi almeno $2^{k-1}a_{2^{k}}$. Gli indici da $2$
a $2^{m}$ sono l'unione di questi blocchi con $k=1,\dots,m$, e
\[
  s_{2^{m}}\ge a_{1}+\sum_{k=1}^{m}2^{k-1}a_{2^{k}}=a_{1}+\tfrac12\,(t_{m}-a_{1})\ge\tfrac12\,t_{m} .
\]
Le due serie sono a termini non negativi, e per il \cref{teo:serie-positive} basta
confrontare la limitatezza delle somme parziali. Se $t_{m}\le T$ per ogni $m$, dato $n$ si
ha $n\le2^{n+1}-1$, e poiché $s_{n}$ è crescente $s_{n}\le s_{2^{n+1}-1}\le t_{n}\le T$: le $s_{n}$ sono limitate.
Se invece $s_{n}\le S$ per ogni $n$, la prima stima dà $t_{m}\le2s_{2^{m}}\le2S$. Dunque le due
serie convergono o divergono insieme.
\end{dimostrazione}

Passando al limite nelle \eqref{eq:condensazione-stime}, la somma della serie data risulta
compresa tra la metà della somma della serie condensata e la somma stessa. Il criterio
richiede soltanto la monotonia, non una funzione continua di cui calcolare l'integrale:
è la versione «discreta» del criterio dell'integrale.

\begin{esempio}[L'armonica generalizzata senza integrali]\label{es:condensazione-armonica}
Ritroviamo il \cref{teo:armonica-generalizzata}. Se $p\le0$, il termine generale non
tende a zero. Se $p>0$, la successione $a_{n}=n^{-p}$ è positiva e decrescente, e
\[
  2^{k}a_{2^{k}}=2^{k}\,(2^{k})^{-p}=\bigl(2^{1-p}\bigr)^{k} :
\]
la serie condensata è la serie geometrica di ragione $2^{1-p}>0$, che converge se e solo
se $2^{1-p}<1$, cioè $p>1$ (\cref{teo:serie-geometrica}). Per $p=1$ la serie condensata
è $1+1+1+\cdots$: è l'argomento di Oresme. Per $p>1$ si ha anche
$\sum_{n\ge1}\frac1{n^{p}}\le\frac1{1-2^{1-p}}$, che per $p=2$ vale $2$.
\end{esempio}

\begin{esempio}[Serie con il logaritmo]\label{es:condensazione-log}
Sia $p>0$. I termini della serie $\sum_{n=2}^{\infty}\frac1{n(\log n)^{p}}$ sono positivi e
decrescenti, perché il denominatore cresce; ponendo $a_{1}=a_{2}$, il che non cambia il carattere (\cref{teo:resto-serie}),
si può applicare il criterio. Per $k\ge1$
\[
  2^{k}a_{2^{k}}=\frac{2^{k}}{2^{k}\,(k\log2)^{p}}=\frac1{(\log2)^{p}}\cdot\frac1{k^{p}} ,
\]
e la serie condensata è, a meno della costante $(\log2)^{-p}$, la serie armonica
generalizzata di esponente $p$. Dunque $\sum_{n\ge2}\frac1{n(\log n)^{p}}$ converge se e solo se
$p>1$. Per $p=1$ si ottiene la divergenza di $\sum\frac1{n\log n}$, una serie che, come
vedremo nel \S\ref{sec:criteri-convergenza}, sfugge al criterio degli infinitesimi.
\end{esempio}

\begin{notastorica}[Oresme, i Bernoulli, Eulero]
La terza dimostrazione della divergenza della serie armonica risale a Nicola d'Oresme,
che intorno al 1350 raggruppò i termini a blocchi di lunghezza $1,2,4,8,\dots$, ciascuno di
somma almeno $\frac12$: sono i blocchi da $\frac1{n+1}$ a $\frac1{2n}$ con $n=2^{j}$: è il criterio di
condensazione applicato alla serie armonica. La dimostrazione cadde nell'oblio e il risultato fu ritrovato da Pietro
Mengoli (1650) e dai fratelli Jakob e Johann Bernoulli (1689). Il problema di calcolare
$\sum\frac1{k^{2}}$, posto da Mengoli, divenne noto come \emph{problema di Basilea}, dalla città
dei Bernoulli, che non riuscirono a risolverlo; lo risolse nel 1734 il giovane Eulero,
che trovò $\frac{\pi^{2}}6$ con un ragionamento audace, reso rigoroso solo in seguito. Nello
stesso periodo Eulero introdusse la costante $\gamma$; Lorenzo Mascheroni, matematico di
Bergamo, ne calcolò nel 1790 trentadue cifre decimali, non tutte esatte.
\end{notastorica}
